\documentclass[10pt]{article}
\usepackage[a4paper,margin=1.7cm]{geometry}
\usepackage[T1]{fontenc}
\usepackage{lmodern}
\usepackage{amsmath,amssymb}
\usepackage{longtable,booktabs,array}
\usepackage{listings}
\usepackage[hidelinks]{hyperref}
\lstset{basicstyle=\ttfamily\footnotesize,breaklines=true,columns=fullflexible,keepspaces=true,upquote=true}
\setcounter{MaxMatrixCols}{20}
\allowdisplaybreaks
\setlength{\parindent}{0pt}\setlength{\parskip}{4pt}\setlength{\emergencystretch}{3em}\sloppy\raggedright
\begin{document}

{\raggedright\LARGE\bfseries Provenance of the README of Pre-gravity6plane\par}\medskip

\begin{itemize}
\item Every equation of Sections 1 to 4 of the README is a restatement of a result computed in 6planes.nb (Wolfram Language 15.0.1, run 2026-10-30, 34 input cells, 0 messages) or verified by verify\_dirac\_matrices.wls (31 exact checks, all True): $T^AT^B+T^BT^A = 2\eta^{AB}I_{16}$; $\Gamma = \mathrm{diag}(-I_8,I_8)$; $\mathcal L^* = \mathcal L$; the Euler-Lagrange equations $T^A\partial_A\psi = (m+\lambda\bar\psi\psi)\psi$ and $\gamma^\mu D_\mu\psi = (m+\lambda\bar\psi\psi)\psi$; $K = -i\sigma_{16}T^4$ with eigenvalues $\pm1$ (8 each); the coupling blocks; $(6p+q)(6p-q) = m^2/H^2$; rank 2 and 28 parameterizations of $M^6$; $\mathcal L_m[\Gamma\psi^*] = \mathcal L_{-m}[\psi]$; $T^0h_{KS}T^0 = -h_{KS}$; the zero energy of pair creation; the bridge data and $R\approx-\frac{31}{216H^2x_0^4}$; the occupation $1/(e^{\epsilon/T_H}+1)$.
\item The matrices $T^A$, $\sigma_{16}$, $\Gamma$, $P_{L,R}$ are those of Pre-gravity6plane.nb (P. L. Nash, (C) 2022, GPL), compared entry by entry with its stored output cells 494, 310, 312, 379, 381.
\item The primordial vielbein is the one displayed in Pre-gravity6plane.nb (text cells 901, 913) with $a_4(Hx_4) = Hx_4$. Definitions of this work: the modern field, with vielbein $\mathrm{diag}(f^{-1/2},1,1,1,f^{1/2},1,1,1)$ and $f = 1-r_s/x_0$; the future field, the same vielbein with $f = 1-r_s/x_0-H_\Lambda^2x_0^2$; and ``Universe'', a mode function (c-number solution) of dirac16complex, whose ``mass'' is its energy for translations of $x_4$ (or the mass parameter $m$ in the mass-sign theorem).
\item Literature: A. Einstein and N. Rosen, Phys. Rev. 48, 73 (1935); W. Kohn and L. J. Sham, Phys. Rev. 140, A1133 (1965); N. D. Mermin, Phys. Rev. 137, A1441 (1965); W. G. Unruh, Phys. Rev. D 14, 870 (1976); J. B. Hartle and S. W. Hawking, Phys. Rev. D 13, 2188 (1976); P. L. Nash, J. Math. Phys. 51, 042501 (2010).
\item Every equation of Section 5 of the README is a result computed in deflation.nb (Wolfram Language 15.0.1, run 2026-10-30, 17 input cells, 0 messages): the line element with $g_{jj} = b_{jj}e^{-ax_4}$; the proper sizes $\sqrt{b_{jj}}e^{-ax_4/2}$; the $x_0$-only dependence of the Ricci scalar, Einstein tensor and Kretschmann scalar; the identity $-dx_4^2 - \sum_j g_{jj}dx_j^2 = -\frac{4}{a^2}(dw^2 + dy_5^2 + dy_6^2 + dy_7^2)/w^2$ with $w = \frac2ae^{ax_4/2}$, $y_j = \sqrt{b_{jj}}x_j$; $R = 3a^2$ and Kretschmann scalar $3a^4/2$ of the pure deflation field; the turning point $x_4^* = \frac1a\ln(b_{55}/p_5^2)$; the field equation of dirac16complex; the dispersion relation for $\omega^2$ and $a_c = 2H$; $M^\dagger K + KM = 0$; and the -mass fraction $1.2\times10^{-31} \to 0.49999995$ of the numerical example ($a = b_{55} = k_1 = m = 1$, $q_5 = 1/2$, $-8 \le x_4 \le 5$).
\item The deflation request and its clarifications are the user's (2026-10-30): $g_{jj} = b_{jj}e^{-ax_4}$ with $a > 0$, $b_{jj} > 0$, entering the line element as $-g_{jj}dx_j^2$, with the primordial field for all other components.
\item Written by the author, P. L. Nash, in the README.md edits of 2026-10-31 (GitHub commits 11ca673 and b6ccd8c): the two sentences on the GNU General Public License that end the first paragraph, the request ``Please cite this work and this web page if you use it.'', and the whole part from the heading Pre-Universe to the end. These are not results of 6planes.nb or deflation.nb. README.tex and README.pdf reproduce them as written, with the author's spelling, punctuation and Mathematica-style notation (for example Subscript[ea, a] and [Dash]).
\item The equations of the Pre-Universe part, as the author wrote them: the quadratic form $x_0^2 + x_1^2 + x_2^2 + x_3^2 - x_4^2 - x_5^2 - x_6^2 - x_7^2$ of $\mathbb R^{4,4}$ (written with x0, ..., x7); $e^{4,4}_Ae^{4,4}_B + e^{4,4}_Be^{4,4}_A = id\ g_{AB}$ for $A, B = 0, 1, \dots, 7$, with $g_{AB} = \mathrm{diag}(1,1,1,1,-1,-1,-1,-1)$; $16 = 2^{8/2}$; $\psi\in\mathbb R^{16} = \{\psi_1\in\mathbb R^8\oplus\psi_2\in\mathbb R^8\}$ with $\psi_1 = \psi_1^a$, $\psi_2 = \psi_2^{(a)}$, $a = 0, 1, \dots, 7$; $\langle xy, xy\rangle = \langle x,x\rangle\langle y,y\rangle$ and $\langle 1,1\rangle = 1$ for the inner product $\langle\ ,\ \rangle: Os\times Os\to\mathbb R$ (written \textless{} , \textgreater{} by the author); $Os = \mathbb R\oplus Im\ Os$; $x^* = 2\langle x,1\rangle - x$; $e_0 = 1$; $x = x^Ae_A$; $\langle e_A, e_B\rangle = g_{A,B}$; $e_Ae_B = e_Cm^C_{AB}$; $\langle u,u\rangle = 1$; $\Psi_1 = \psi_1^ae^{1/2}_a$; $\Psi_2 = \psi_2^{(a)}e^{1/2}_{(a)}$; and the split octonion multiplication table (row $A$, column $B$: $e_Ae_B$)
\end{itemize}

\[
\begin{array}{ccccccccc}
\text{A/B} & e_0 & e_1 & e_2 & e_3 & e_4 & e_5 & e_6 & e_7 \\
e_0 & e_0 & e_1 & e_2 & e_3 & e_4 & e_5 & e_6 & e_7 \\
e_1 & e_1 & -e_0 & -e_3 & e_2 & -e_5 & e_4 & e_7 & -e_6 \\
e_2 & e_2 & e_3 & -e_0 & -e_1 & -e_6 & -e_7 & e_4 & e_5 \\
e_3 & e_3 & -e_2 & e_1 & -e_0 & -e_7 & e_6 & -e_5 & e_4 \\
e_4 & e_4 & e_5 & e_6 & e_7 & e_0 & e_1 & e_2 & e_3 \\
e_5 & e_5 & -e_4 & e_7 & -e_6 & -e_1 & e_0 & -e_3 & e_2 \\
e_6 & e_6 & -e_7 & -e_4 & e_5 & -e_2 & e_3 & e_0 & -e_1 \\
e_7 & e_7 & e_6 & -e_5 & -e_4 & -e_3 & -e_2 & e_1 & e_0
\end{array}
\]

\begin{itemize}
\item Checked for this record (Python, 2026-10-31) from the table above, with $\langle x,y\rangle = g_{AB}x^Ay^B$: $e_0$ is a two-sided unit; $e_Ae_A = -g_{AA}e_0$ for $A = 1, \dots, 7$; $e_Ae_B = -e_Be_A$ for $A \ne B$, $A, B \ge 1$; $\langle xy, xy\rangle = \langle x,x\rangle\langle y,y\rangle$ holds for 20000 random pairs of integer vectors $x$, $y$ with components in $-9..9$, and $x(xy) = (xx)y$ for 5000 pairs with components in $-5..5$.
\item Works cited in the Pre-Universe part: R. L. Bryant, ``Submanifolds and special structures on the octonians'', J. Differential Geometry 17 (1982) 185-232; A. Einstein and N. Rosen, ``The Particle Problem in the General Theory of Relativity'', Phys. Rev. 48, 73-77 (1935), doi:10.1103/PhysRev.48.73; P. A. M. Dirac, ``A remarkable representation of the 3 + 2 de Sitter group'', J. Math. Phys. 4, 901-909 (1963); E. A. Lord, ``The Dirac spinor in six dimensions'', Proc. Camb. Phil. Soc. 64, 765-778 (1968).
\item README.tex is made from README.md by code/md2tex.py, which reads the Markdown as GitHub does: the HTML comments are not shown (README.tex keeps them as LaTeX comments); \texttt{\textbackslash{}[Dash]}, \texttt{\textbackslash{}[CapitalPsi]} and the other backslash escapes show as [Dash], [CapitalPsi]; in math, \texttt{\textbackslash{}\textbackslash{}\{} and \texttt{\textbackslash{}\textbackslash{}\}} are braces and the MathJax macro \texttt{\textbackslash{}space} is a space; a \texttt{\$\$...\$\$} that is a whole paragraph is displayed, and one inside a paragraph is set inline in display style; formulas written in GitHub's verbatim math syntax (a code block marked math, or inline TeX between dollar-backtick delimiters) are set as display and inline math; the indented Dynkin diagram is set verbatim; the line in the span with font-size 30px is set at 18.75pt, that is 30px relative to GitHub's 16px text, scaled to the 10pt text of README.pdf (GitHub itself drops the style attribute and shows the line at normal size); text and tables are set ragged right, as on GitHub.
\item On 2026-10-31 the formulas of Sections 1 to 5 of the README that GitHub displayed wrongly were rewritten in that verbatim syntax: GitHub's Markdown processing had turned \texttt{\textbackslash{}\{} into \texttt{\{}, \texttt{\textbackslash{},} into a comma and the row break \texttt{\textbackslash{}\textbackslash{}} of a matrix into a space, and had paired two \texttt{*} as emphasis; the label $\Delta\mathrm{SCF}$ became one formula. The TeX of these formulas is unchanged.
\end{itemize}

\end{document}